\scanpage{119}
\uline{Cor} Dans le cas de la prop 2, \deleted{la catégorie} on a
\[
G^{\unsure{\circ|0}}\simeq U^{\unsure{\circ|0}}\times G_m,\quad\varepsilon=\operatorname{pr}_2,\quad j(\lambda)=\deleted{\unsure{\operatorname{pr}_2(\lambda)^2}}\operatorname{inj}_{\unsure{1|2}}(\lambda)^2.
\]
La structure \deleted{motivique 1) 2) 3)} de catégorie tannakienne de $M_{\mathbb R}$ est connue totalement par la connaissance de $G$ et l'élém de $H^2(\mathbb R,\unsure{U})$ caractéristique, qui \deleted{\unsure{provient}} \deleted{\unsure{est}} l'unique élém non nul de $H^2(\mathbb R,U)$, correspondant à l'unique élément non nul de $H^2(\mathbb R,\mu_2)$. Il existe un \deleted{élém} objet $E$ de $M_{\mathbb R}$ \deleted{\unsure{de}} de degré 1, rang 2, \deleted{\unsure{non trivial}} tel que la sous-catégorie tannakienne \deleted{$\unsure{M'}$} $M'$ qu'il engendre sera déjà $\mathbb R$-tordue, et \deleted{\unsure{correspond à un quotient}} \deleted{\unsure{qui}} \unsure{par} $U^{\unsure{\circ|0}}$ ``y opère \uline{trivialt}''.

Dans le cas favorable $\mathscr G_{\mathbb R}\ne\varnothing$, \margin{\inserted{i.e. $U$ un tore} \deleted{\unsure{à degrés pairs}}} l'ens \deleted{$\operatorname{Ob}\mathscr G_{\mathbb R}$} \inserted{$H^0(\mathscr G_{\mathbb R})$ des objets à isom près de $\mathscr G_{\mathbb R}$} est un \deleted{ens principal homogène} \inserted{torseur} sous le groupe $H^1(\mathbb R,U)$, qui lui-même est isom. can. à ${}_2U(\mathbb R)\underset{\text{non can.}}{\simeq}(\mathbb Z/2\mathbb Z)^{\dim U}$. D'autre part, à chaque élém de $\operatorname{Ob}\mathscr G_{\mathbb R}$ correspond une \deleted{structure de \unsure{fonctorialité}} \inserted{structure polarisante} sur $G^{\unsure{\circ|0}}$, définie par un hom
\diagram{\begin{tikzcd}
\mu_{2\mathbb C}\arrow[r,"i_1"] & G^{\unsure{\circ|0}}_{\mathbb C}
\end{tikzcd}}
tel que $\varepsilon_{\mathbb R}i_1=\mathrm{id}$, $i_1\bar i_1=j_{\mathbb R}$, et l'ens de ces structures est égalemt un torseur sous ${}_2U(\mathbb R)$. L'application
\diagram{\begin{tikzcd}
H^0(\mathscr G_{\mathbb R})\arrow[r] & \operatorname{Pol}(\mathscr G_{\mathbb R})
\end{tikzcd}}
est compatible avec l'opération de ${}_2U(\mathbb R)$, \uline{donc bijective}. Donc les $\mathbb R$-foncteurs fibres à isom
\mspage{118}
